\subsection{Diagnostic Classification and MCI Stage Discrimination}
Classification used independent subject-level three-fold evaluation for AD versus CN ($n=470$), EMCI+LMCI versus CN ($n=741$), and EMCI versus LMCI ($n=436$). The MCI/CN task includes EMCI and LMCI, while MCI-U is excluded. The completed analysis comprises 63 imaging and 18 clinical fits (27 task--input configurations, three folds each). Results are reported as the equally weighted fold mean $\pm$ sample s.d., alongside conditional bootstrap 95\% CIs (Table~\ref{tab:classification_summary}). This is cross-sectional MCI stage discrimination, not future conversion or longitudinal prognostic prediction.

For AD/CN, MRI+Syn-FDG+Syn-AV45 achieved ACC $88.08\pm2.06\%$ and AUC $93.82\pm3.21\%$, compared with MRI-only ACC $79.56\pm6.20\%$ and AUC $87.02\pm5.44\%$ (Fig.~\ref{fig:5}A). The paired AUC difference was $+6.80$ percentage points (95\% CI $[+3.39,+10.22]$). For MCI/CN, MRI+Syn-both had ACC $65.45\pm1.64\%$ and AUC $69.46\pm1.34\%$, versus MRI-only $67.48\pm0.23\%$ and $71.59\pm1.80\%$; the paired AUC difference was $-2.13$ points (95\% CI $[-6.30,+2.10]$; Fig.~\ref{fig:5}B). Thus the synthetic combination did not improve all tasks.

For EMCI/LMCI, MRI+Syn-both had ACC $65.82\pm2.06\%$ and AUC $66.33\pm0.54\%$, compared with MRI-only $64.90\pm3.09\%$ and $62.81\pm2.50\%$. Its balanced-ACC difference favored the synthetic combination by $6.91$ points (95\% CI $[2.02,11.98]$), while the ACC and AUC difference CIs crossed zero (Fig.~\ref{fig:5}C). The Syn-minus-Real dual-PET AUC intervals crossed zero for all three tasks, which does not establish equivalence or interchangeability (Table~\ref{tab:classification_comparisons}).

Clinical-only AUC was $100.00\pm0.00\%$ for AD/CN, $98.93\pm0.70\%$ for MCI/CN, and $99.94\pm0.08\%$ for EMCI/LMCI. MRI+Clinical AUC was $99.17\pm0.96\%$, $93.07\pm6.10\%$, and $89.60\pm9.08\%$, respectively. MRI+Syn-both AUC was below both clinical baselines in each task, with the specified paired intervals below zero. These observations do not support independent incremental diagnostic value over the clinical variables, which include diagnosis-related scores and diagnosis-conditioned fills. The zero observed AD/CN error for Clinical-only and its degenerate empirical bootstrap CI do not establish perfect population generalization.

\begin{figure}[htbp]
\centering
\includegraphics[width=\textwidth]{figure_5}
\caption{\textbf{Cross-sectional classification under the revised three-fold protocol.} Panels A--C correspond to AD/CN, EMCI+LMCI/CN and EMCI/LMCI. Each panel compares nine input conditions: MRI; MRI+Real FDG; MRI+Syn FDG; MRI+Real AV45; MRI+Syn AV45; MRI+Real-both; MRI+Syn-both; Clinical-only; and MRI+Clinical. Values are equally weighted three-fold means; error bars show sample s.d. in percentage points, not 95\% CIs. ACC, AUC, SEN, SPE and balanced ACC are reported with their CIs in Table~\ref{tab:classification_summary}. The results do not demonstrate longitudinal prognosis, PET equivalence or benefit independent of diagnosis-related clinical inputs.}
\label{fig:5}
\end{figure}

Supplemental ROC curves use fold-wise averaging and pointwise confidence bands (Fig.~\ref{fig:classification_roc}). Paired DeLong tests are reported by outer fold with their raw and family-specific Holm-adjusted $p$ values (Table~\ref{tab:classification_delong}); they are not combined into a three-fold overall DeLong test. All comparisons are conditional on the fixed trained models and splits.
